%%%PAGE 273 rot=0 legibility=good summary=机械能守恒(链条非质点、守恒条件辨析、连接体型、阻力正比于速度下落的功率能量图像)、板块连接体极限法、读数坑与游标卡尺分度、动量定理时间
%%%BLOCK id=273-01 ch=06 sec=机械能守恒·链条 kind=derivation alt=none cont=none y=0.05-0.21
%FIG p273_a 0.085 0.063 0.212 0.205 soft
\notefig[0.25]{figures/p273_a.png}

非质点\unclear{系}

由机械能守恒
\[
\frac12 mg\frac L2=\frac12 mv^2
\]
\[
v=\sqrt{\frac{gL}{2}}
\]
%%%END
%%%BLOCK id=273-02 ch=06 sec=机械能守恒·条件辨析 kind=knowledge alt=none cont=none y=0.05-0.19
\begin{itemize}
\item 合外力为0 不是\textsuperscript{机守}守恒守恒 % CHECK: 此行“守恒”似连写两遍；“机守”为写在上方的插入小字
\item 合外力做功为0 也不是机械能守恒\quad（匀速放\unclear{重}物） % 括号内为写在行末右下方的小字
\item 只有重力、弹力做功\\
  $\to$ 系统内的弹力
\end{itemize}
%%%END
%%%BLOCK id=273-03 ch=06 sec=机械能守恒·连接体型 kind=method alt=04 cont=none y=0.22-0.35
\textbf{机械能守恒}

连接体型
\begin{itemize}
\item 速度关联不同
\item $h_1$、$h_2$ 高度变化不同
\end{itemize}
%%%END
%%%BLOCK id=273-04 ch=06 sec=功能关系·图像 kind=diagram alt=03 cont=none y=0.38-0.53
% 四幅图自左至右：v-t；P_{E_p}/P_G-t；E_K-t(下方标 P_{合外力})；E-t(下方标 P_f)
$f=kv$ 下落

%FIG p273_b 0.07 0.385 0.82 0.525
\notefig[0.85]{figures/p273_b.png}
%%%END
%%%BLOCK id=273-05 ch=03 sec=连接体·极限法 kind=problem alt=none cont=none y=0.54-0.67
%FIG p273_c 0.06 0.53 0.32 0.645
\notefig[0.3]{figures/p273_c.png}

极限法
\[
T=3.5+M\,\frac{mg-3.5}{1+M}
\]% CHECK: 手写分母为“1+M”(疑为 1+m)；图中 M=1kg、m=0.5kg、摩擦因数标注为0.35
\[
m\to+\infty\qquad T=13.5\unit{N}
\]
%%%END
%%%BLOCK id=273-06 ch=18 sec=读数·易错 kind=note alt=none cont=none y=0.72-0.85
\textbf{读数坑}

%FIG p273_d 0.03 0.745 0.25 0.815
\notefig[0.3]{figures/p273_d.png}

$\Delta x=x_2-2x_1$
%%%END
%%%BLOCK id=273-07 ch=18 sec=游标卡尺·读数 kind=note alt=none cont=none y=0.86-0.89
游标卡尺副尺几个格就是几分度 % 原稿“副尺”为插入的小字，写在“尺”“几”之间的上方
%%%END
%%%BLOCK id=273-08 ch=07 sec=动量定理 kind=note alt=none cont=none y=0.89-0.93
动量定理时间，注意题目问的是全程时间\unclear{哪}还是阶段时间 % “还”前的字被涂写重叠，疑为“哪”
%%%END
%%%PAGEEND 273
